\section{Tessuto cartilagineo}

Il \textbf{tessuto connettivo di sostegno} fornisce impalcatura resistente per l'inserzione delle formazioni muscolari, per assicurare la pervietà di alcuni organi cavi (trachea e laringe), per la protezione di visceri e per il sostegno di tutto l'organismo. Comprende:
\begin{itemize}
  \item tessuto cartilagineo;
  \item tessuto osseo.
\end{itemize}

\rubrica{Distribuzione della cartilagine}
Nell'embrione forma lo scheletro primitivo, che alla nascita viene sostituito da tessuto osseo. La \textit{cartilagine metafisaria} è presente nelle ossa lunghe per tutto il periodo di accrescimento. Nell'adulto permangono:
\begin{itemize}
  \item cartilagini articolari;
  \item cartilagini intercostali;
  \item cartilagine nasale;
  \item cartilagini laringee;
  \item anelli tracheali;
  \item cartilagini bronchiali;
  \item padiglione auricolare;
  \item anello fibroso (dischi intervertebrali).
\end{itemize}

% ---------------------------------------------------------------------------
\subsection{Cellule}

Le cellule, dette \textbf{condrociti}, sono scarse e sono poste in gruppi di 2--4 cellule (\textbf{gruppi isogeni}) in piccoli spazi all'interno della sostanza fondamentale, dette \textbf{lacune}.

\rubrica{Linea differenziativa}
Cellula mesenchimale $\rightarrow$ cellula condroprogenitrice $\rightarrow$ condroblasto $\rightarrow$ condrocita.

\rubrica{Modalità di accrescimento del tessuto cartilagineo}
\begin{itemize}
  \item \textbf{accrescimento per apposizione}: in superficie, a partire dallo strato condrogenico interno del pericondrio, che produce nuova matrice cartilaginea;
  \item \textbf{accrescimento interstiziale}: a carico dei condrociti già presenti all'interno del tessuto cartilagineo, che si dividono formando i gruppi isogeni.
\end{itemize}
Il pericondrio è costituito da uno strato fibroso esterno (fibrociti e collagene di tipo I) e da uno strato condrogenico interno.

% ---------------------------------------------------------------------------
\subsection{Caratteristiche}

Costituito in prevalenza da sostanza fondamentale: un gel compatto che contiene derivati polisaccaridici detti \textbf{condroitinsolfati}, complessati con proteine a formare \textbf{proteoglicani}.

\rubrica{Funzioni}
\begin{itemize}
  \item fornisce sostegno rigido, ma più flessibile dell'osso;
  \item riduce l'attrito tra le superfici ossee articolari;
  \item resiste alla compressione;
  \item limita gli spostamenti reciproci di ossa vicine.
\end{itemize}

\rubrica{Proprietà meccaniche}
I proteoglicani (PGA), paragonabili a una spugna molecolare, si circondano di molecole d'acqua, che attraggono con i loro gruppi caricati negativamente presenti in superficie. In presenza di un carico o di una sollecitazione meccanica, l'acqua, essendo indeformabile, abbandona la sostanza amorfa, lasciando interagire i gruppi caricati negativamente dei proteoglicani che, per repulsione elettrostatica, si oppongono al carico pressorio. PGA tipici del tessuto cartilagineo sono l'aggrecano e il versicano.

% ---------------------------------------------------------------------------
\subsection{Tipi di cartilagine}

Le fibre sono presenti in tipi e quantità diverse a seconda del tipo di cartilagine:

\begin{itemize}
  \item \textbf{cartilagine ialina}: fibre collagene immerse in sostanza fondamentale, liscia; cartilagini articolari e di sostegno delle vie respiratorie;
  \item \textbf{cartilagine elastica}: fibre elastiche, leggera e flessibile; padiglione auricolare ed epiglottide laringea;
  \item \textbf{cartilagine fibrosa}: numerose fibre collagene e scarsa sostanza fondamentale, dura e resistente; dischi intervertebrali.
\end{itemize}

% ---------------------------------------------------------------------------
\subsection{Cartilagine ialina}

Traslucida, trasparente, di colore bianco-azzurrognolo.
\begin{itemize}
  \item componente amorfa ricca di glicoproteine e proteoglicani;
  \item componente fibrillare costituita da collagene;
  \item condrociti disposti in gruppi isogeni di 2--4 cellule;
  \item non vascolarizzata;
  \item delimitata da connettivo fibroso con funzione trofica e condrogenica $\rightarrow$ \textbf{pericondrio}.
\end{itemize}

\rubrica{Localizzazione}
Cartilagini articolari, cartilagini intercostali, cartilagine nasale, cartilagini laringee, anelli tracheali, cartilagini bronchiali.

\rubrica{Componente cellulare}
Aggregati di cellule mesenchimatiche diventano rotondeggianti $\rightarrow$ \textbf{centri di condrificazione}. Si differenziano in \textbf{condroblasti}: cellule secernenti che producono matrice cartilaginea (proteoglicani e collagene), la quale, interponendosi fra i condroblasti stessi, ne determina l'allontanamento. Le cellule restano imprigionate all'interno di lacune scavate nella matrice stessa $\rightarrow$ \textbf{condrociti}: cellule quiescenti. I condroblasti possono dividersi all'interno della lacuna formando gruppi di 2--4 cellule $\rightarrow$ \textbf{gruppi isogeni} o cloni cellulari.

\rubrica{Matrice extracellulare}
\begin{itemize}
  \item \textbf{componente amorfa}: acqua (80\%); sali minerali, tra cui NaCl; proteoglicani (30--40\%); glicoproteine;
  \item \textbf{componente fibrillare}: collagene (40\%), con diametro 15--30\,nm nella zona periferica (rete tridimensionale lassa) e 45\,nm nella zona centrale.
\end{itemize}

% ---------------------------------------------------------------------------
\subsection{Cartilagine articolare}

Priva di pericondrio; la funzione trofica è svolta dal liquido sinoviale. Consente alle articolazioni di sopportare pressioni meccaniche da varie direzioni.
\begin{itemize}
  \item \textbf{strato tangenziale}: cellule allungate;
  \item \textbf{strato intermedio}: cellule tondeggianti;
  \item \textbf{strato radiale}: gruppi isogeni paralleli tra loro, perpendicolari alla superficie.
\end{itemize}

% ---------------------------------------------------------------------------
\subsection{Cartilagine metafisaria}

Presente nelle ossa lunghe durante l'accrescimento, si organizza in zone concentriche, dalla periferia verso il centro dell'osso:
\begin{itemize}
  \item \textbf{CR}: cartilagine a riposo;
  \item \textbf{CS}: cartilagine seriata (in proliferazione);
  \item \textbf{CI}: cartilagine ipertrofica;
  \item \textbf{CD}: cartilagine in degenerazione (calcificazione della matrice intercellulare, elevata mineralizzazione);
  \item \textbf{ZI}: zona di invasione e deposizione del tessuto osseo.
\end{itemize}
Nella zona di invasione si distinguono matrice ossea (O), matrice cartilaginea (C) e midollo osseo (M).

% ---------------------------------------------------------------------------
\subsection{Cartilagine elastica}

Presenza di abbondanti fibre elastiche associate a fibre collagene: elastica e flessibile.
\begin{itemize}
  \item le cellule presentano un grosso vacuolo $\rightarrow$ nucleo e citoplasma spinti alla periferia $\rightarrow$ forma ad anello con castone;
  \item gruppi isogeni meno numerosi, con meno condrociti;
  \item le fibre elastiche formano una rete tridimensionale, più spesse e abbondanti nella porzione centrale, più sottili nella regione periferica.
\end{itemize}

\rubrica{Localizzazione}
Padiglione auricolare, epiglottide, tromba di Eustachio, laringe.

% ---------------------------------------------------------------------------
\subsection{Cartilagine fibrosa}

Presenza di abbondanti fibre collagene e scarsa componente amorfa; colore biancastro. Somigliante a un tessuto connettivo denso regolare.
\begin{itemize}
  \item cellule rotondeggianti, contenute in una lacuna in cui sono presenti proteoglicani;
  \item cellule isolate, raramente raggruppate in gruppi isogeni disposti in file parallele.
\end{itemize}

\rubrica{Localizzazione}
Dischi intervertebrali (anello fibroso).
